%%%%%%%%%%%%%%%%%%%%%%%%%%%%%%%%%%%%%%%%%%%%%%%%%%
%%%%  TEMPLATE IEEE PARA ENTREGA DE INFORMES  %%%% 
%%%% DE LABORATORIO DE LA UNIVERSIDAD CENTRAL %%%%
%%%%%%%%%%%%%    BOGOTÁ, COLOMBIA    %%%%%%%%%%%%%
%%%%%%%%%%%%%%%%%%%%%%%%%%%%%%%%%%%%%%%%%%%%%%%%%%
%%%%   AUTOR: SERGIO ANDRÉS CHAPARRO MORENO   %%%%
%%%%%%%%%%%%%%%%%%%%%%%%%%%%%%%%%%%%%%%%%%%%%%%%%%
%%%%%%%%%%%%%  VERSIÓN 1.1-FEB 2017  %%%%%%%%%%%%%
%%%%%%%%%%%%%%%%%%%%%%%%%%%%%%%%%%%%%%%%%%%%%%%%%%
\documentclass[journal]{IEEEtran}
\IEEEoverridecommandlockouts
%%%%%%%%%%%%%%%%%%%%%%%%%%%%%%%%%%%%%%
%%%%%%%% PRINCIPALES PAQUETES %%%%%%%%
%%%%%%%%%%%%%%%%%%%%%%%%%%%%%%%%%%%%%%
\usepackage{fancyhdr}
\usepackage{graphicx}
\usepackage[spanish, es-tabla]{babel}
\usepackage[utf8]{inputenc}
\usepackage{color}
\usepackage{hyperref}
\usepackage{wrapfig}
\usepackage{array}
\usepackage{multirow}
\usepackage{adjustbox}
\usepackage{nccmath}
\usepackage{subfigure}
\usepackage{amsfonts,latexsym} % para tener disponibilidad de diversos simbolos
\usepackage{enumerate}
\usepackage{booktabs}
\usepackage{float}
\usepackage{threeparttable}
\usepackage{array,colortbl}
\usepackage{ifpdf}
\usepackage{rotating}
\usepackage{cite}
\usepackage{stfloats}
\usepackage{url}
\usepackage{listings}
%%%%%%%%%%%%%%%%%%%%%%%%%%%%%%%%%%%%%%%%%%%
%%% CREAR Y REESCRIBIR ALGUNOS COMANDOS %%%
%%%%%%%%%%%%%%%%%%%%%%%%%%%%%%%%%%%%%%%%%%%
\newcolumntype{P}[1]{>{\centering\arraybackslash}p{#1}}  %% Se crea un nuevo tipo de columna llamada P.
\newcommand{\tabitem}{~~\llap{\textbullet}~~}
\newcommand{\ctt}{\centering\scriptsize\textbf} %%\ctt abrevia el comando \centering\scriptsize\textbf
\newcommand{\dtt}{\scriptsize\textbf} %%\dtt abrevia el comando \scriptsize\textbf
\renewcommand\IEEEkeywordsname{Palabras clave}
%%%%%%%%%%%%%%%%%%%%%%%%%%%%%%%%%%%%%%%%%%%


% correct bad hyphenation here
\hyphenation{op-tical net-works semi-conduc-tor} %% Con este comando se especifican como pueden seprarse las sílabas adecuadamente en caso una palabra quede en dos lineas diferentes de texto

\graphicspath{ {figs/} }  %%Ruta donde se encuentran las imágenes, que esté vacio indica que las imagenes están dentro de la misma carpeta que contiene el archivo .tex


%%%%%%%%%%%%%%%%%%%%%%%%%%%%%%%%%%%%%%%%%%%%%%%%%%%%%%%%%%
%%% ENCABEZADO DE LAS PÁGINAS TIPO UNIVERSIDAD CENTRAL %%%
%%%%%%%%%%%%%%%%%%%%%%%%%%%%%%%%%%%%%%%%%%%%%%%%%%%%%%%%%%
\newcommand{\MYhead}{\smash{\scriptsize
\hfil\parbox[t][\height][t]{\textwidth}{\centering
\begin{picture}(0,0) \put(-0,-15){\includegraphics[width=25mm]{LogoUCentral}} \end{picture} \hspace{6.4cm}
INFORME DE PRÁCTICA DE LABORATORIO \hspace{5.15cm} Versión 1.1\\
\hspace{6.2cm} DEPARTAMENTO DE INGENIERÍA ELECTRÓNICA \hspace{4.45cm} Periodo 2017-1\\
\underline{\hspace{ \textwidth}}}\hfil\hbox{}}}
\makeatletter
% normal pages
\def\ps@headings{%
\def\@oddhead{\MYhead}%
\def\@evenhead{\MYhead}}%
% title page
\def\ps@IEEEtitlepagestyle{%
\def\@oddhead{\MYhead}%
\def\@evenhead{\MYhead}}%
\makeatother
% make changes take effect
\pagestyle{headings}
% adjust as needed
\addtolength{\footskip}{0\baselineskip}
\addtolength{\textheight}{-1\baselineskip}
%%%%%%%%%%%%%%%%%%%%%%%%%%%%%%%%%%%%%%%%%%%%%%%%%%%%%%%%%%
%%%%%%%%%%%%%%%%%%%%%%%%%%%%%%%%
%%%%% INICIO DEL DOCUMENTO %%%%%
%%%%%%%%%%%%%%%%%%%%%%%%%%%%%%%%
\begin{document}
%%%%%%%%%%%%%%%%%%%%%%%%%%%%
%%% TÍTULO DEL DOCUMENTO %%%
%%%%%%%%%%%%%%%%%%%%%%%%%%%%
\title{Título de la práctica}
%%%%%%%%%%%%%%%%%%%%%%%%%%%%
%%%%%%%%% AUTORES %%%%%%%%%
%%%%%%%%%%%%%%%%%%%%%%%%%%%
\author{Apellido1~Nombre1,~
        Apellido2~Nombre2,~
        y~Apellido3~Nombre3\\
				\textit{\{Estudiante1,~Estudiante2~y~Estudiante3\}@ucentral.edu.co}\\
				Profesor:~Apellido~Nombre\\% stops a space
\thanks{El presente documento corresponde a un informe de práctica de laboratorio de Electrónica III presentado en la Universidad Central durante el periodo 2017-1.}} %\thanks anexa una nota a pie de página donde se puede colocar alguna información sobre la naturaleza del documento.
%%%%%%%%%%%%%%%%%%%%%%%%%%%

% Comando que indica la generación del título
\maketitle

%%%%%%%%%%%%%%%%%%%%%
%%%%%% RESUMEN %%%%%%
%%%%%%%%%%%%%%%%%%%%%
\begin{abstract}
En este documento se presenta una plantilla elaborada en \LaTeX  para la presentación de informes de laboratorio del departamento de ingeniería electrónica de la universidad central. En esta sección se espera el resumen de la propuesta de laboratorio seguido de los experimentos realizados y los principales resultados alcanzados.
\end{abstract}
% En el resumen no se recomienda colocar citaciones bibliográficas.
%%%%%%%%%%%%%%%%%%%%%%
%%% PALABRAS CLAVE %%%
%%%%%%%%%%%%%%%%%%%%%%
\begin{IEEEkeywords}
IEEE, plantilla, \LaTeX, ecuaciones, figuras, referencias, tablas. Escribir las palabras clave que se encuentran en el documento.
\end{IEEEkeywords}
%%%%%%%%%%%%%%%%%%%%%%
%\IEEEpeerreviewmaketitle
%%%%%%%%%%%%%%%%%%%%%%%%%%%%%%%%%%%%%
%%% PRIMERA SECCIÓN DEL DOCUMENTO %%%
%%%%%%%%%%%%%%%%%%%%%%%%%%%%%%%%%%%%%
\section{Introducción}
\IEEEPARstart{E}{scribir} la propuesta de laboratorio y un pequeño marco teórico.

\subsection{Marco teórico}
Presentar brevemente los diferentes conceptos necesarios para el desarrollo de la práctica. Esto puede hacer en varias subsub-secciones.

\subsubsection{Concepto1}
Este es un ejemplo de subsubsección.

\subsubsection{Concepto2}

%%%%%%%%%%%%%%%%%%%%%%%%%%%%%%%%%%%%%
%%% SEGUNDA SECCIÓN DEL DOCUMENTO %%%
%%%%%%%%%%%%%%%%%%%%%%%%%%%%%%%%%%%%%
\section{Trabajo previo}
Resolver las preguntas previas a la práctica de laboratorio.

%%%%%%%%%%%%%%%%%%%%%%%%%%%%%%%%%%%%
%%% TERCERA SECCIÓN DEL DOCUMENTO %%%
%%%%%%%%%%%%%%%%%%%%%%%%%%%%%%%%%%%%
\section{Desarrollo de la práctica}
En esta sección se describen los diferentes experimentos, montajes, cálculos o simulaciones hechas durante la práctica de laboratorio, junto a los resultados obtenidos representados mediante gráficas y tablas. Cada punto de la práctica puede numerarse o colocarse como una subsección con un título adecuado.

\subsection{Figuras en \LaTeX}
Para anexar una gráfica de datos se recomienda que sea en formato .eps o .ps lo cual puede hacerse usando MATLAB como se muestra en \cite{imagenes}. 
%%%%%%%%%%%%%%%%%%%%%%%%%%%%%%%%%%%%%%
%%%%%%%%  ANEXAR UNA GRÁFICA  %%%%%%%%
%%%%%%%%%%%%%%%%%%%%%%%%%%%%%%%%%%%%%%
\begin{figure}[H] %[H] obliga a la figura a quedar en la misma posición en el texto final que en el archivo .tex, [t] coloca la figura en la parte superior de la página, [b] coloca la figura en la parte inferior de la página. 
\centering  %Centra la figura
\includegraphics[scale=0.23]{fig} %[se define tamaño de la figura]{nombre del archivo con la figura}
\caption{Nombre descriptivo de la figura.} %Numera y titula la gráfica
\label{lvdt4} %Permite referenciar la grafica en el texto EJ: en la gráfica \ref{lvdt4} se observa...
\end{figure}
También se pueden anexar subfiguras, modificar la posición y el tamaño. En el archivo de imagen no debe haber título. Si se desea anexar imágenes extraídas de otras fuentes (por ejemplo Internet), estas deben tener buena calidad y preferiblemente estar en formato png).

Para referenciar o nombrar una figura en el texto: En la figura \ref{lvdt4} se presenta la característica $I_1$ contra $V_1$.

Por otro lado, para crear esquemáticos de circuitos o diagramas de bloques o de flujo, puede usarse el software \emph{DIA} \cite{dia}, u otros programas que permitan salvar preferiblemente gráficos en formatos .eps o .ps. 
\begin{figure}[H] %[H] obliga a la figura a quedar en la misma posición en el texto final que en el archivo .tex, [t] coloca la figura en la parte superior de la página, [b] coloca la figura en la parte inferior de la página. 
\centering  %Centra la figura
\includegraphics[scale=0.55]{LVDT4} %[se define tamaño de la figura]{nombre del archivo con la figura}
\caption{Diagrama del integrado AD598.} %Numera y titula la gráfica
\label{diafig} %Permite referenciar la grafica en el texto EJ: en la gráfica \ref{lvdt4} se observa...
\end{figure}
%%%%%%%%%%%%%%%%%%%%%%%%%%%%%%%%%%%%%%

%%%%%%%%%%%%%%%%%%%%%%%%%%%%%%%
%%%%%%%%% ECUACIONES %%%%%%%%%%
%%%%%%%%%%%%%%%%%%%%%%%%%%%%%%%
\subsection{Ecuaciones en \LaTeX}
Para escribir una ecuación:
\begin{equation}\label{eqID}
I_D=\frac{q N_A n_i^2}{N_D}\left(\frac{\alpha V_{GS}^2}{\mu_o}\right)^3
\end{equation}

\begin{equation}\label{Voeq} %label se usa para referenciar la ecuación al igual que con las gráficas.
V_o \approx \int e^XdX
\end{equation}
Para mencionar una ecuación en el texto: en (\ref{eqID}) y (\ref{Voeq}) se observan las relaciones para $I_D$ y $V_o$ respectivamente.

Se pueden reportar despejes, cálculos y procedimientos sin enumerarlos. Por ejemplo el siguiente cálculo:
\begin{gather*}
i=\frac{v}{R}\Longrightarrow i=\frac{5}{500}=10 mA
\end{gather*}
%%%%%%%%%%%%%%%%%%%%%%%%%%%%%%%

%%%%%%%%%%%%%%%%%%%%%%%%%%%
%%%%%%%%% TABLAS %%%%%%%%%%
%%%%%%%%%%%%%%%%%%%%%%%%%%%
\subsection{Tablas en \LaTeX}
Para definir una tabla:

\begin{table}[H]
\centering
\caption{Nombre de la tabla}
\label{table1}
\begin{tabular}{c c c}\hline\hline
\textbf{Símbolo} & \textbf{Nombre} & \textbf{Código Latex}\\ \hline
$\alpha$ & alpha & \verb|\alpha| \\
$\mu$ & mu & \verb|\mu|\\
$\beta$ & beta & \verb|\beta|\\
$\Omega$ & Omega & \verb|\Omega| \\\hline \hline
\end{tabular}
\end{table}
Para mencionar la tabla en el documento: en la tabla \ref{table1} se muestran algunos ejemplos de código \LaTeX para obtener letras griegas.
%%%%%%%%%%%%%%%%%%%%%%%%%%%

%%%%%%%%%%%%%%%%%%%%%%%%%%%%%%
%%%%% CITAR BIBLIOGRAFIA %%%%%
%%%%%%%%%%%%%%%%%%%%%%%%%%%%%%
\subsection{Citar en formato IEEE}
Para citar referencias bibliográficas se usa el comando cite. En \cite{nombre_para_citar} se muestran los campos que deben llenarse en una referencia, en \cite{kopka} se muestra un ejemplo, y en \cite{link} se muestra como citar un enlace. Preferiblemente citar libros y artículos.
%%%%%%%%%%%%%%%%%%%%%%%%%%%%%%

%%%%%%%%%%%%%%%%%%%%%%%%%%%%%%%%%%%%
%%% CUARTA SECCIÓN DEL DOCUMENTO %%%
%%%%%%%%%%%%%%%%%%%%%%%%%%%%%%%%%%%%
\section{Conclusiones}
Reportar en tercera persona las diferentes conclusiones producto de la práctica de laboratorio desarrollada.
%%%%%%%%%%%%%%%%%%%%%%%%%%%%%%%%%%%%

\ifCLASSOPTIONcaptionsoff
  \newpage
\fi

%%%%%%%%%%%%%%%%%%%%%%%%%%
%%%%%% BIBLIOGRAFIA %%%%%%
%%%%%%%%%%%%%%%%%%%%%%%%%%
\begin{thebibliography}{1}

\bibitem{imagenes}
Youtube, canal schaparro. \url{https://youtu.be/IhvF6iY7n5k}. Recuperado el 30 de Enero de 2017.

\bibitem{dia}
Dia Diagram Editor. \url{https://sourceforge.net/projects/dia-installer/}. Recuperado el 30 de Enero de 2017.

\bibitem{nombre_para_citar}
Inicial1.~Apellido1 and Inicial2.~Apellido2, \emph{Nombre de libro}, \#edición~ed.\hskip 1em plus
  0.5em minus 0.4em\relax Ciudad, País: Editorial, año.

\bibitem{kopka}
H.~Kopka and P.~W. Daly, \emph{A Guide to \LaTeX}, 3rd~ed.\hskip 1em plus
  0.5em minus 0.4em\relax Harlow, England: Addison-Wesley, 1999.

\bibitem{link}
Overleaf. \url{https://www.overleaf.com/}. Recuperado el 02 de Febrero de 2017.

\end{thebibliography}
%%%%%%%%%%%%%%%%%%%%%%%%%%

\end{document}
%%%%%%%%%%%%%%%%%%%%%%%%%%%%%%%%
%%%%%% FIN DEL DOCUMENTO %%%%%%%
%%%%%%%%%%%%%%%%%%%%%%%%%%%%%%%%



